\section{The question and the analytical structure}\label{sec:intro}

The Green approach asks for a signed upper estimate on an observable
whose zero representation is positive. Its usefulness rests on keeping
these two descriptions of the same quantity together. Positivity of a
source, regularity of its heat flow and accuracy of an approximation
are valuable inputs, but none by itself estimates that signed value.

This manuscript continues the Green-potential investigation and integrates
the relevant foundations from both of its parent manuscripts
\cite{GNPBackground,GNPOperator,OGParent}. It is organized around the chain
\[
 \text{complete theta source}\ \longrightarrow\ H_\lambda
 \ \longrightarrow\ \log|H_\lambda|\ \longrightarrow\ J_\lambda
 \ \longrightarrow\ \text{a quantitative defect comparison}.
\]
The spatial variable is \(z\), Newman time is \(\lambda\), and forward
heat time is \(\tau=\lambda_0-\lambda\). The half-line cosine convention
is fixed in \eqref{eq:H}; in particular \(H_0(z)=\xi((1-\ii z)/2)/8\).
The original source variable and the observed spatial variable are
distinct. A heat operation on one need not induce the same operation on
the other.

Sections~\ref{sec:source}--\ref{sec:log-collision} derive the source,
geometric and operator dictionaries and explain why collision singularities
belong to root and logarithmic coordinates while the source orbit stays
regular. Sections~\ref{sec:green}--\ref{sec:prime} identify the complete
Green potential and its exact time-zero arithmetic baseline. The boundary
logarithm, finite-source and deflation arguments in
Sections~\ref{sec:boundary-log}--\ref{sec:deflation} give distinct ways of
acquiring that observable. Section~\ref{sec:RS} explains how classical and
heat-evolved Riemann--Siegel formulas enter those estimates.
Sections~\ref{sec:transition}--\ref{sec:euler-reference} then develop
the child's results through E06: the uniform outer reduction, the
limits of quadratic acquisition, and the Euler-reference comparison.
Theorem~\ref{thm:b10-averaged} extends the residual mean to a growing
block-constant prime cutoff on shorter translated intervals at \(c=1\).
Corollary~\ref{cor:b10-green} gives the corresponding Green mean and
fixed-threshold exceptional proportion. Each deduction keeps the
physical time fixed inside the spatial integral at each center.

The conditioning example in Section~\ref{sec:conditioning} concerns a
specified class containing noncanonical translated sources. The
short-mollifier theorem in Section~\ref{sec:quadratic} concerns one
polynomial shared by a whole interval of centers. Their hypotheses and
quantifiers are part of the conclusions. Neither is a general obstruction
to theta-specific signed arithmetic estimates. Section~\ref{sec:parallel}
keeps localized correlations and collision constraints as parallel
investigations, identifying exactly which inputs they could supply.

There are four different levels of assertion throughout this paper:
an exact representation, an estimate of acquisition error, a conditional
algorithm with unpaid data requirements, and an upper value small enough
for an exclusion. The first three do not imply the fourth. The one-defect
threshold \eqref{eq:margin} provides a numerical scale against which an
estimate can be judged; a reformulation of the all-height target is
identified once at its first use. The premise screen in
Section~\ref{sec:evidence} applies the repository's collision controls
\cite{OGControls}; the source/evidence map accompanying this manuscript
points to the result and avoidance indexes and the exact inherited files.

The author-facing purpose is a coherent entry point for further Green
research. The child has selected a transition-scale consumer and proved
both acquisition estimates and scoped obstructions. Its outstanding
issue is the arithmetic value at exceptional centers, rather than
selection of another representation. The final section distinguishes
that issue from the parent's time-zero question and from the additional
uniformity needed for other defect depths or time bands. This synthesis
does not assert a new theorem of canonical smallness.
